\subsection*{Operation Type}

Croqtile maps these functionalities at the software level as different \textbf{DMA Operation Types}. The supported operations include:

\begin{itemize}[leftmargin=1.6em]
\item \texttt{dma.copy}: Copies flat memory directly.
\item \texttt{dma.pad}: Pads shaped data while transferring the data.
\item \texttt{dma.transp}: Transposes shaped data while transferring the data.
\end{itemize}

The configurations, such as padding and transposing details, are programmed as parameters of the DMA operation. For example:

\begin{lstlisting}[language=croqtile,caption={Operation Type}]
global f32 [32, 16, 9] input;
dma.transp<0, 2, 1> input => shared;  // Result shape [32, 9, 16]
dma.pad<{1, 0, 3}, {0, 1, 2}, {0, 0, 0}, 0.1f> input => shared; // Result shape [33, 17, 14]
\end{lstlisting}

Here, the \textbf{DMA configuration}s are enclosed by \texttt{\textless{}\textgreater{}}. The configuration varies according to different operations. The detailed configuration syntax and limitations for a specific platform, such as , are listed below as an example:

Programmers should note that support for DMA types other than \texttt{dma.copy} varies by platform. For instance, on the  platform, it could be either a mapping of TMA or an orchestration of load instructions from multiple threads.
